\printedpage{49}
% head: A. Grothendieck : Dérivateurs, Ch. VI, Édition des Universités de Montpellier II et Paris VII
% foot: Version 9 novembre 2001

\[
\begin{tikzcd}
\mathcal{F}'\Sigma^{-1} \arrow[r, "\approx"] & \mathcal{F}\Sigma^{-1} \arrow[l, bend left=40, "\approx"', "\Psi_{\mathcal{F}',\mathcal{F}}"]
\end{tikzcd}
\qquad\qquad (\Sigma = W_S \text{ or } W_S^{\mathrm{g}})
\]

for two full subcategories $\mathcal{F}'$, $\mathcal{F}$
\[
\usl{Cofib}\,S \subseteq \mathcal{F}' \subseteq \mathcal{F} \subseteq \mathrm{Cat}/S.
\]
The functors $\Phi$, $\Psi$ send the full subcategory $\usl{Fibhot}\,S$ of ``the upper floor'', into the subcategories $\usl{Fibsc}_0S$, $\usl{Cofibsc}_0S$, respectively, hence induce on the upper floor functors $\Phi_0$, $\Psi_0$
\[
\usl{Fibsc}_0S \xleftarrow{\Phi_0} \usl{Fibhot}\,S \xrightarrow{\Psi_0} \usl{Cofibsc}_0S,
\]
playing on the upper floor exactly the same role as $\Phi$, $\Psi$ on the ground floor. The only difference is that on the upper floor we have $W_S = W_S^{\mathrm{d}} = W_S^{\mathrm{g}}$ $(= W_S^{\mathrm{f}})$,
\msp{page 78}
so there is only one sequence of localized categories to examine. We thus find
\[
\begin{tikzcd}
(\usl{Fibsc}_0S)W_S^{-1} \arrow[r, "\approx"] & \mathcal{F}W_S^{-1} \arrow[l, bend left=40, "\approx"', "\Phi_{0,\mathcal{F}}"]
\end{tikzcd}
\]
\[
\begin{tikzcd}
\mathcal{F}'W_S^{-1} \arrow[r, "\approx"] & \mathcal{F}W_S^{-1} \arrow[l, bend left=40, "\approx"', "\Phi_{0,\mathcal{F}',\mathcal{F}}"]
\end{tikzcd}
\]
for $\mathcal{F}$ (or $\mathcal{F}$, $\mathcal{F}'$) full subcategory(ies) of $\usl{Fibhot}\,S$ such that
\[
\usl{Fib}_0S \subseteq \mathcal{F} \subseteq \usl{Fibhot}\,S \qquad\text{resp.}\qquad \usl{Fib}_0S \subseteq \mathcal{F}' \subseteq \mathcal{F} \subseteq \usl{Fibhot}\,S,
\]
and likewise on the right-hand side for the full subcategories $\mathcal{F}$ or $\mathcal{F}$, $\mathcal{F}'$ of $\usl{Fibhot}\,S$ satisfying
\[
\usl{Cofib}_0S \subseteq \mathcal{F} \subseteq \usl{Fibhot}\,S \qquad\text{resp.}\qquad \usl{Cofib}_0S \subseteq \mathcal{F}' \subseteq \mathcal{F} \subseteq \usl{Fibhot}\,S,
\]
NB. The category $\usl{Parf}\,S$ contains both $\usl{Fib}_0S$ and $\usl{Cofib}_0S$, and is therefore among the admissible categories just like $\mathcal{F}$, or [?] $\mathcal{F}'$.

Thus, the four functors $\Phi$, $\Psi$, $\Phi_0$, $\Psi_0$ give all the quasi-inverses one can wish for on the ground floor on the one hand ($\Phi$ and $\Psi$), on the upper floor on the other ($\Phi_0$, $\Psi_0$). They therefore reduce all the relevant localized categories to the following four
\[
\begin{gathered}
(\mathrm{Cat}/S)W_S^{-1},\ \ (\mathrm{Cat}/S)(W_S^{\mathrm{d}})^{-1},\ \ (\mathrm{Cat}/S)(W_S^{\mathrm{g}})^{-1} \qquad \text{(on the ground floor)}\\
(\usl{Fibhot}\,S)W_S^{-1} \qquad \text{(on the upper floor).}
\end{gathered}
\]

\encadre{VII} It remains to compare the localized category on the upper floor with the localized categories on the ground floor. The vertical arrows descending from the upper floor to the ground floor, are compatible with $W_S^{\mathrm{d}}$ on the left-hand side of the diagram, with $W_S^{\mathrm{g}}$ on the right-hand side, so they induce essentially \textit{two} different functors
\printedpage{50}
% head: A. Grothendieck : Dérivateurs, Ch. VI, Édition des Universités de Montpellier II et Paris VII
% foot: Version 9 novembre 2001
\msp{page 79}
$i$, $j$ on $(\usl{Fibhot}\,S)W_S^{-1}$, into $(\mathrm{Cat}/S)(W_S^{\mathrm{d}})^{-1}$ and $(\mathrm{Cat}/S)(W_S^{\mathrm{g}})^{-1}$:
\[
\begin{tikzcd}[column sep=small, row sep=large]
& (\usl{Fibhot}\,S)W_S^{-1} \arrow[dl, "i"'] \arrow[dr, "j"] \arrow[dd, dotted, "\approx"', "\delta"] & \\
\mathrm{HOT}_W(S) \overset{\text{def}}{\approx} (\mathrm{Cat}/S)(W_S^{\mathrm{d}})^{-1} \arrow[dr, dotted, "\rho"'] & & (\mathrm{Cat}/S)(W_S^{\mathrm{g}})^{-1} \overset{\text{def}}{\approx} \mathrm{HOT}^{o}_{W}(S) \arrow[dl, dotted, "\sigma"] \\
& \begin{array}{l} (\mathrm{Cat}/S)W_S^{-1} \\ \overset{\text{def}}{\approx} \mathrm{HOT}^{\mathrm{lc}}_{W}(S) \end{array} &
\end{tikzcd}
\qquad ,
\]
so as to make the square completed by the dotted lines commute, representing the localization functors $\rho$, $\sigma$. We have of course
\[
\begin{aligned}
(*) \qquad \rho i = \sigma j : (\usl{Fibhot}\,S)W_S^{-1} \longrightarrow (\mathrm{Cat}/S)W_S^{-1}\\
\text{induced by the inclusion}\\
\delta : \usl{Fibhot}\,S \hookrightarrow \mathrm{Cat}/S
\end{aligned}
\]
We have stated in \encadre{I} that this arrow is an equivalence, so that $\rho$ and $\sigma$ can be interpreted, up to isomorphism of functors, as retractions of $\mathrm{HOT}_W(S)$ and of $\mathrm{HOT}^{o}_{W}(S)$ onto a common subcategory, interpreted either as $(\usl{Fibhot}\,S)W_S^{-1}$, or as $(\mathrm{Cat}/S)W_S^{-1}$. It remains to describe the quasi-inverse functor of $(*)$, which will at the same time give the quasi-inverses of all the vertical functors (from the upper floor to the ground floor) in the diagram p. 76, more precisely for their localized versions, for $W_S$.

\msp{page 80}
This is what is obtained with the help of the ``ascending'' functor (from the ground floor to the upper floor) $C$, which one can regard as a functor $\mathrm{Cat}/S \longrightarrow \usl{Fibhot}\,S$, but which in fact takes its values in $\usl{Parf}\,S$. (This is a detail of no great importance, it seems to me.)
\[
\mathrm{Cat}/S \xrightarrow{C} \usl{Parf}\,S \hookrightarrow \usl{Fibhot}\,S.
\]
This functor is compatible with the localizers $W_S$ on both sides, and establishes equivalences (quasi-inverses of the equivalences
\[
\begin{gathered}
(\usl{Parf}\,S)W_S^{-1} \xrightarrow{\sim} (\mathrm{Cat}/S)W_S^{-1} \qquad\text{and}\qquad (\usl{Fibhot}\,S)W_S^{-1} \xrightarrow{\sim} (\mathrm{Cat}/S)W_S^{-1}),\\
(\mathrm{Cat}/S)W_S^{-1} \xrightarrow{\sim} (\usl{Parf}\,S)W_S^{-1} \left( \xrightarrow{\sim} (\usl{Fibhot}\,S)W_S^{-1} \right).
\end{gathered}
\]

\printedpage{51}
% head: A. Grothendieck : Dérivateurs, Ch. VI, Édition des Universités de Montpellier II et Paris VII
% foot: Version 9 novembre 2001
So, when the three HOT categories for $S$ are interpreted in terms of localized categories of $\mathrm{Cat}/S$, it is the functor $C$ (more precisely the composite $\mathrm{Cat}/S \xrightarrow{C} \usl{Parf}\,S \hookrightarrow \mathrm{Cat}/S$) that provides both the inclusion functor (!) $i$, and the inclusion functor $j$ into $\mathrm{HOT}_W(S)$ and $\mathrm{HOT}^{o}_{W}(S)$ respectively. The functor in question
\[
C' : \mathrm{Cat}/S \longrightarrow \mathrm{Cat}/S
\]
does indeed satisfy
\[
C'(W_S) \subseteq W_S^{\mathrm{d}} \cap W_S^{\mathrm{g}}
\]
and defines the two functors
\[
\begin{tikzcd}[column sep=large]
& (\mathrm{Cat}/S)(W_S^{\mathrm{d}})^{-1} \approx \mathrm{HOT}_W(S) \\
\begin{array}{l} (\mathrm{Cat}/S)W_S^{-1} \\ \approx \mathrm{HOT}^{\mathrm{lc}}_{W}(S) \end{array} \arrow[ur, "i"] \arrow[dr, "j"'] & \\
& (\mathrm{Cat}/S)(W_S^{\mathrm{g}})^{-1} \approx \mathrm{HOT}^{o}_{W}(S)
\end{tikzcd}
\]
This allows us to make explicit the adjunction formulas of \encadre{IV}, when we interpret $\xi$, $\eta$ there
\msp{page 81}
as coming from $X$, $Y$ in $\mathrm{Cat}/S$:
\[
\boxed{\mathrm{Hom}_{W_S^{\mathrm{d}}}(X,CY) \simeq \mathrm{Hom}_{W_S^{\mathrm{g}}}(X,CY) \simeq \mathrm{Hom}_{W_S}(X,Y),}
\]
where $\mathrm{Hom}_{W_S^{\mathrm{d}}}$, $\mathrm{Hom}_{W_S^{\mathrm{g}}}$ and $\mathrm{Hom}_{W_S}$ denote respectively the Hom in the localized categories for $W_S^{\mathrm{d}}$, $W_S^{\mathrm{g}}$ and $W_S$, i.e. in $\mathrm{HOT}_W(S)$, $\mathrm{HOT}^{o}_{W}(S)$ and $\mathrm{HOT}^{\mathrm{lc}}_{W}(S)$ respectively (\footnote[51]{This formula suggests that if $X \in \mathrm{Ob}\,\mathrm{Cat}/S$, $Y \in \mathrm{Ob}\,\usl{Fibhot}\,S$, then
\[
\mathrm{Hom}_{W_S^{\mathrm{d}}}(X,Y) \simeq \mathrm{Hom}_{W_S^{\mathrm{g}}}(X,Y) \simeq \mathrm{Hom}_{W_S}(X,Y),
\]
but this follows from the adjunction formula and from $Y \longrightarrow C(Y)$ in $W_S^{\mathrm{d}} \cap W_S^{\mathrm{g}}$ for $Y$ in $\usl{Fibhot}\,S$.}). % [sic] adjoinction -> adjonction, spelling (elsewhere "adjonction")

\textbf{Comment.} All the constructions and statements from \encadre{I} to \encadre{VI}, with the exception of the fact that the descending functors induce equivalences for the \textit{coarse} localized categories, and of the adjunction formulas of \encadre{IV}, are of a virtually sorital nature, and require no conditions on the basic localizer $W \subseteq \mathrm{Fl}\,\mathrm{Cat}$ other than L1, L3 quater, L4 bis, summarized as a) b) c) on page 45. (We do not even have to use d).) It is true that in the absence of condition L5 (page 43), we do not know whether on the upper floor the two relevant localizers, $W_S$ and $W_S^{\mathrm{f}}$ (equal to $W_S^{\mathrm{d}}$ on the left-hand side, to $W_S^{\mathrm{g}}$ on the right-hand side) are equal. Thus a priori we find on the upper floor \textit{two} essentially distinct localized categories, and not just one, namely $(\usl{Fibhot}\,S)\Sigma^{-1}$ with $\Sigma = W_S$ or $\Sigma = W_S^{\mathrm{f}}$ $(= W_S^{\mathrm{d}} = W_S^{\mathrm{g}})$.

\printedpage{52}
% head: A. Grothendieck : Dérivateurs, Ch. VI, Édition des Universités de Montpellier II et Paris VII
% foot: Version 9 novembre 2001
On the other hand, the fact that the descent functors induce equivalences between the coarse localized categories is deeper, it requires (it seems to me) that we know an ``ascending'' functor $C$, say from $\mathrm{Cat}/S$ to $\usl{Fibhot}\,S$, compatible with $W_S$, and equipped with
\[
X \xrightarrow{i_X} C(X),
\]
which is in $W_S$ for arbitrary $X$. When we know that Cat is a Quillen category, even without having the properties of the functors $\usl{Ch}_\infty$ or $\usl{Ch}$ (\footnote[52]{nor even being sure of the existence of a \textit{functorial} factorization of $X \xrightarrow{f} S$ into $X \xrightarrow{i_X} C(X) \xrightarrow{\bar{f}} S$ with $i_X \in W$ and $\bar{f}$ perfect.}), we should probably
\msp{page 82}
be able to manage. But the adjunction formula seems to me to remain unexpected and deep, even admitting the results of \textsc{Quillen} and \textsc{Thomason} (\footnote[53]{See however theorem-scholie 2, p. 108-113. It is after all not as deep as it seemed to me, but virtually sorital!}). It seems linked here to a rather special property of the functor $C$, namely the lemma p. 68
\[
\begin{gathered}
i_{C(X)} \underset{\big\uparrow}{=} C(i_X) : C(X) \longrightarrow C(C(X)),\\
\scriptstyle \text{in } \underbrace{(\mathrm{Cat}/S)(W_S^{\mathrm{d}})^{-1}}_{\mathrm{HOT}_W(S)} \text{ and in } \underbrace{(\mathrm{Cat}/S)(W_S^{\mathrm{g}})^{-1}}_{\mathrm{HOT}^{o}_{W}(S)}
\end{gathered}
\]
which does not look purely formal. It would be interesting to examine to what extent this formula still makes sense in a more or less arbitrary ``category of models'' $(\mathcal{M},W)$, or coming from a Quillen triple $(W,\mathrm{Cof},\mathrm{Fib})$ in $\mathcal{M}$. It is clear that the factorization $X \xrightarrow{f} S$ into $X \xrightarrow{i_X} C(X) \xrightarrow{\bar{f}} S$ $(i_X \in W,\ \bar{f} \in \mathrm{Parf})$, playing the role of ``fibrations'' in the Quillen context, has an intrinsic meaning independent of the choice of a particular construction, something we can make perfectly precise. It is the description of the localizers $W_S^{\mathrm{d}}$, $W_S^{\mathrm{g}}$ in $\mathcal{M}/S$ (for $S \in \mathrm{Ob}\,\mathcal{M}$) that is not understood in a general way, outside the case $\mathcal{M} = \mathrm{Cat}$. In other words, I do not understand how to construct categories $\mathrm{HOT}_{\mathcal{M},W}(S)$, $\mathrm{HOT}^{o}_{\mathcal{M},W}(S)$, even if $(\mathcal{M},W)$ comes from a Quillen triple, I see only the description of a category $\mathrm{HOT}^{\mathrm{lc}}_{\mathcal{M},W}(S) = (\mathcal{M}/S)(W_S^{\text{coarse}})^{-1}$, and a localization functor $\mathcal{M}/S \longrightarrow \mathrm{HOT}^{\mathrm{lc}}_{\mathcal{M},W}(S)$. We have to manage to describe in $\mathcal{M}/S$ $W$-equivalences (\textit{``relative to $S$''}) finer than the coarse one given by $W$ on the source $X$ and nothing more. In Cat I manage it by using either the family of local isomorphisms (to describe $W_S^{\mathrm{d}}$), or that of colocal isomorphisms (to describe $W_S^{\mathrm{g}}$) - in any case, there is a notion of ``localization'' (either on $S$, or on $S^{o}$ \dots), with % [sic] En d’autre termes -> En d’autres termes, fixed expression requires plural
\msp{page 83}
all the richness of the intuitions linked to such a notion. Perhaps, without knowing it I have always treated Cat not as a ``bare'' category, but as a site, or rather (since there
\printedpage{53}
% head: A. Grothendieck : Dérivateurs, Ch. VI, Édition des Universités de Montpellier II et Paris VII
% foot: Version 9 novembre 2001
is the involution $X \longmapsto X^{o}$ which cannot be evaded), as a ``bi-site'' - we have \textit{two} ways of localizing there, either by the $X/x$ (Zariski topology!) or by the $x\backslash X$ (co-Zariski topology). There is also the smooth topology (probably very close to the intuition of étale cohomology), and the proper topology - and we expect the smooth one to give essentially the same results as the Zariski one (since the étale topology is equal to it - there is no need to make a fuss about it \dots); and the proper topology the same results (e.g. the same cohomology?) as the co-Zariski topology. % [sic] vue que -> vu que, conjunctional locution

It does seem that in all the categories of models used up to now to describe Hot as a localized category, we have notions of localization: (Top), $\mathbf{\Delta}^\wedge$ (or $A^\wedge$, with $A$ an arbitrary test category), (Cat); and in Cat we even have \textit{two} such notions. On the other hand I am thinking of $\infty$-groupoids, it is perhaps the only category of models already so close to Hot itself, that we discern no notion of localization in it, and that the very idea of wanting to localize in it (except by ``localizing'' Hot itself) seems almost
\msp{page 84}
preposterous.

This quick digression now convinces only that the said compatibility for the construction of $C(X)$ and $i_X : X \longrightarrow C(X)$, can in no case be purely formal, a mere matter of Quillen categories say. That will at best give the compatibility in $(\mathcal{M}/S)W_S^{-1}$, which is ridiculous! But if we take as localizer $\Sigma$ in $\mathcal{M}/S$ the isomorphisms and nothing more, the compatibility relation will become false every time! As a slightly more reasonable set of arrows $\Sigma$, which is still very small, we can think e.g. of \textit{universal} $W$-equivalence over $S$ (distinctly finer still than $W_S^{\mathrm{d}}$, $W_S^{\mathrm{g}}$ in the special case of Cat). Would we still have compatibility for this set of arrows? I am also a little perplexed by the fact (not well understood) that the adjunction formulas for the Zariski and co-Zariski topologies go in opposite directions. (Or might I have made a mistake in dualizing in a hurry?) But to tell the truth, this compatibility by itself, even when we use a site structure on $\mathcal{M}$, is not enough to write down an adjunction formula. We have two localized categories, % [sic] ne convainc -> me convainc, the verb needs an object pronoun
\[
\mathrm{HOT}^{\mathrm{lc}}_{\mathcal{M}} = (\mathcal{M}/S)W_S^{-1}, \qquad \mathrm{HOT}_W(S) = (\mathcal{M}/S)(W_S^{\mathrm{f}})^{-1},
\]
where $W_S^{\mathrm{f}}$ denotes a ``fine'' $W$-equivalence, corresponding to the notion of localization, more precisely to the ``admissible'' base changes we have chosen, and we have a localization functor
\msp{page 85}
\[
\mathrm{Hot}_W(S) \longrightarrow \mathrm{Hot}^{\mathrm{lc}}_{W}(S),
\]
but no functor in the reverse direction. But yes, but yes, working with the similar category $\usl{Fibhot}\,S$, this must still make sense at least when we have a good notion of ``fibration''. I see a little work of conceptual roughing-out to do, with perhaps an adjunction formula at the end of it, but I am not going to launch into it now.